% Lösungen Einheit 5
\einheitenkopf[l5][5 Anwendungen]{Einheit 5 von 5 · Anwendungen}

\erg{45}{Anfangswert – Änderung je Einheit: a) 4\,€ – 2,20\,€ je km \quad b) 80\,l – 5\,l je Minute (weniger) \quad c) 25\,€ – 9,99\,€ je Monat \quad d) 45\,€ – 3\,€ je Woche \quad e) 12\,cm – 0,5\,cm je Tag}
\erg{46}{a) $y = 2x + 6$ \quad b) $y = 3x + 15$ \quad c) $y = 0{,}3x + 12$ \quad d) $y = 0{,}05x + 2$ (5 Cent = 0,05\,€)}
\erg{47}{a) $y = 4x + 8$ \quad b) 20\,€ \quad c) $y = -15x + 250$ \quad d) 130\,€ \quad e) $y = 35x + 60$; $35x + 60 = 500$ ergibt $x \approx 12{,}6$; nach 13 Monaten (nach 12 Monaten erst 480\,€, nach 13 Monaten 515\,€)}
\erg{48}{zum Beispiel: a) Verleih: 10\,€ Grundgebühr, 2,50\,€ je Stunde \quad b) Tank mit 60\,l, je Minute laufen 3\,l ab}
\erg{49}{a) II \quad b) I \quad c) III: beginnt bei 20\,€ auf der $y$-Achse und steigt je Besuch um 5\,€ (flacher als II)}
\erg{50}{a) Nord: $9 + 0{,}3 \cdot 40 = 21$\,€, Süd: $0{,}45 \cdot 40 = 18$\,€; Süd ist günstiger \quad b) Nord: 33\,€, Süd: 36\,€; Nord ist günstiger \quad c) Nord: $y = 0{,}3x + 9$, Süd: $y = 0{,}45x$; $0{,}3x + 9 = 0{,}45x$, $x = 60$; ab 60\,km ist Nord günstiger}
\erg{51}{a) 2\,€ \quad b) $2 + 30 \cdot 0{,}25 = 9{,}50$\,€ \quad c) Plus: $5 + 10 \cdot 0{,}15 = 6{,}50$\,€; Plus ist günstiger}
\erg{52}{a) Emma hat Cent und Euro gemischt; richtig $y = 0{,}2x + 1{,}5$ \quad b) Paul hat Anfangswert und Änderung vertauscht; richtig $y = 45x + 30$}
\erg{53}{a) Ja: Die Geraden sind parallel, die erste liegt immer um den Unterschied der Grundgebühren darunter. \quad b) Nein: Die erste Gerade steigt steiler; ab einer bestimmten Strecke schneiden sich die Geraden, danach ist der zweite Tarif günstiger.}
